\documentclass[../main.tex]{subfiles}
\graphicspath{{\subfix{../img/}}{\subfix{../architectures images/}}}
\begin{document}
\chapter{Sprint 5 : Boutique et Gestion de Stock}

%──────────────────────────────────────────────
\section{Introduction}
%──────────────────────────────────────────────

Beaucoup de salles vendent des compl\'{e}ments, des v\^{e}tements ou des accessoires sur place. Ce sprint construit la cha\^{i}ne compl\`{e}te : fournisseur $\rightarrow$ achat $\rightarrow$ stock par salle $\rightarrow$ vente au membre. Les transferts inter-salles permettent aux partenaires multi-centres de r\'{e}\'{e}quilibrer leur inventaire sans friction. Le tout est prot\'{e}g\'{e} par le \texttt{ModuleGuard} --- la boutique n'est accessible qu'aux partenaires qui l'ont souscrite.

%──────────────────────────────────────────────
\section{Sprint Backlog}
%──────────────────────────────────────────────

Le tableau~\ref{tab:backlog_sprint5} pr\'{e}sente le backlog du Sprint 5 d\'{e}taill\'{e} par \'{e}pic, nature et user stories.

\begin{longtable}{|p{2.3cm}|p{2.8cm}|p{3.7cm}|p{5.2cm}|}
\hline
\textbf{Epic} & \textbf{Nature} & \textbf{User Story} & \textbf{Description} \\
\hline
\endfirsthead
\hline
\textbf{Epic} & \textbf{Nature} & \textbf{User Story} & \textbf{Description} \\
\hline
\endhead
\hline
\endfoot

\multirow{2}{*}{\parbox{2.5cm}{Gestion des produits}} & \acrshort{crud} & En tant qu'administrateur, je veux g\'{e}rer les produits & Cr\'{e}er, modifier, consulter et supprimer des produits avec leurs informations : nom, prix unitaire, association au partenaire et image. \\
\cline{2-4}
 & Consultation & En tant qu'administrateur, je veux consulter le catalogue & Visualiser l'ensemble des produits disponibles avec filtres et recherche. \\
\hline

\multirow{2}{*}{\parbox{2.5cm}{Gestion des fournisseurs}} & \acrshort{crud} & En tant qu'administrateur, je veux g\'{e}rer les fournisseurs & Cr\'{e}er et g\'{e}rer les fiches fournisseurs pour le suivi des approvisionnements. \\
\cline{2-4}
 & Consultation & En tant qu'administrateur, je veux consulter la liste des fournisseurs & Afficher les fournisseurs avec leurs informations de contact et historique d'achats. \\
\hline

\multirow{2}{*}{\parbox{2.5cm}{Achats fournisseurs}} & \acrshort{crud} & En tant qu'administrateur, je veux enregistrer des achats fournisseurs & Saisir les achats avec quantit\'{e}, prix et date de r\'{e}ception pour alimenter le stock. \\
\cline{2-4}
 & Fonctionnalit\'{e} & En tant qu'administrateur, je veux que le stock soit mis \`{a} jour automatiquement & Le stock de la salle concern\'{e}e est incr\'{e}ment\'{e} automatiquement lors de la r\'{e}ception d'un achat. \\
\hline

\multirow{2}{*}{\parbox{2.5cm}{Stock par salle}} & Consultation & En tant qu'administrateur, je veux consulter l'\'{e}tat du stock par salle & Visualiser les quantit\'{e}s disponibles de chaque produit pour chaque salle de sport. \\
\cline{2-4}
 & Fonctionnalit\'{e} & En tant qu'administrateur, je veux recevoir des alertes de stock bas & \^{E}tre notifi\'{e} lorsque le stock d'un produit descend en dessous d'un seuil critique. \\
\hline

\multirow{1}{*}{\parbox{2.5cm}{Transferts de stock}} & Fonctionnalit\'{e} & En tant qu'administrateur, je veux transf\'{e}rer du stock entre salles & Effectuer des mouvements de stock d'une salle source (\texttt{fromGym}) vers une salle destination (\texttt{toGym}) en sp\'{e}cifiant le produit et la quantit\'{e}. \\
\hline

\multirow{2}{*}{\parbox{2.5cm}{Ventes aux membres}} & Fonctionnalit\'{e} & En tant que caissier, je veux vendre des produits aux membres & Enregistrer une vente en s\'{e}lectionnant le produit, la quantit\'{e}, le prix et le membre acheteur. \\
\cline{2-4}
 & Fonctionnalit\'{e} & En tant qu'administrateur, je veux consulter l'historique des ventes & Visualiser l'ensemble des ventes r\'{e}alis\'{e}es avec filtres par date, produit et membre. \\
\hline

\multirow{1}{*}{\parbox{2.5cm}{Feature flag \acrshort{saas}}} & Technique & En tant que syst\`{e}me, je veux contr\^{o}ler l'acc\`{e}s au module boutique & Prot\'{e}ger l'acc\`{e}s au module boutique via le \texttt{ModuleGuard} et le \texttt{SaasModuleEnum.Boutique} selon l'abonnement du partenaire. \\
\hline
\caption{Sprint Backlog -- Sprint 5}
\label{tab:backlog_sprint5}
\end{longtable}

%──────────────────────────────────────────────
\section{Analyse}
%──────────────────────────────────────────────

\subsection{Diagramme de cas d'utilisation}

La figure~\ref{fig:sprint5-usecase} r\'{e}sume les droits par acteur : l'administrateur g\`{e}re le cycle achat/stock/transfert, le caissier effectue les ventes au quotidien.

\begin{figure}[H]
\centering
\begin{tikzpicture}[scale=0.85, transform shape]

% System box
\node[systembox={Syst\`{e}me TrustGym -- Boutique et Stock}, minimum width=10cm, minimum height=12cm] (sys) at (5.5, -5) {};

% Actors (all on the left)
\stickfigure{(-2.5, -2)}{Admin}
\stickfigure{(-2.5, -7)}{Caissier}

% Use cases
\node[usecase] (uc1) at (5.5, 0) {G\'{e}rer les produits};
\node[usecase] (uc2) at (5.5, -1.8) {G\'{e}rer les fournisseurs};
\node[usecase] (uc3) at (5.5, -3.6) {Enregistrer un achat};
\node[usecase] (uc4) at (5.5, -5.4) {Consulter l'\'{e}tat du stock};
\node[usecase] (uc5) at (5.5, -7.2) {Transf\'{e}rer du stock};
\node[usecase] (uc6) at (5.5, -9) {Vendre un produit};
\node[usecase] (uc7) at (5.5, -10.8) {Consulter l'historique des ventes};

% Admin links
\draw (-1, -2) -- (uc1.west);
\draw (-1, -2) -- (uc2.west);
\draw (-1, -2) -- (uc3.west);
\draw (-1, -2) -- (uc4.west);
\draw (-1, -2) -- (uc5.west);
\draw (-1, -2) -- (uc7.west);

% Caissier links
\draw (-1, -7) -- (uc6.west);
\draw (-1, -7) -- (uc4.west);

\end{tikzpicture}
\caption{Diagramme de cas d'utilisation -- Sprint 5}
\label{fig:sprint5-usecase}
\end{figure}

\subsection{Cas d'utilisation cibl\'{e}s}

Format d\'{e}fini en section~\ref{sec:uc_format}. Ce sprint couvre : \textbf{G\'{e}rer les produits} (Admin --- prot\'{e}g\'{e} par \texttt{ModuleGuard}/\texttt{SaasModuleEnum.Boutique}), \textbf{Transf\'{e}rer du stock entre salles} (Admin --- contr\^{o}le de stock source, transaction \texttt{BEGIN}/\texttt{COMMIT}, salles distinctes oblig\'{e}es), et \textbf{Vendre un produit \`{a} un membre} (Caissier/Admin --- d\'{e}cr\'{e}mentation atomique du \texttt{ProductStock}).

%──────────────────────────────────────────────
\section{R\'{e}alisation}
%──────────────────────────────────────────────

\subsection{Backend}

\begin{longtable}{|>{\footnotesize}p{4.2cm}|>{\footnotesize}p{3.6cm}|p{5.9cm}|}
\hline
\textbf{Module} & \textbf{Entit\'{e}s cl\'{e}s} & \textbf{R\^{o}le} \\
\hline
\endfirsthead
\hline
\endfoot
\texttt{ProductModule} & \texttt{Product}, \texttt{ProductCategory} & \acrshort{crud} produits et cat\'{e}gories ; filtrage par partenaire. \\
\hline
\texttt{SupplierModule} & \texttt{Supplier}, \texttt{SupplierPurchase} & Fiches fournisseurs et achats ; mise \`{a} jour automatique du stock \`{a} la r\'{e}ception. \\
\hline
\texttt{ProductStockModule} & \texttt{ProductStock} & Inventaire par salle (produit $\times$ salle $\times$ quantit\'{e}). \\
\hline
\texttt{StockTransferModule} & \texttt{StockTransfer} & Transfert inter-salles transactionnel (\texttt{fromGym}/\texttt{toGym}) avec v\'{e}rification de stock. \\
\hline
\texttt{MemberPurchaseModule} & \texttt{MemberPurchase} & Ventes aux membres ; d\'{e}cr\'{e}mentation atomique du \texttt{ProductStock}. \\
\hline
\texttt{StockMovementModule} & \texttt{StockMovement} & Journal d'audit complet : chaque entr\'{e}e/sortie horodat\'{e}e avec raison et responsable. \\
\hline
\caption{Modules backend du Sprint 5}
\end{longtable}

\subsection{Diagramme de s\'{e}quence}

La figure~\ref{fig:sprint5-sequence} montre que le transfert s'ex\'{e}cute en transaction : soit les deux mises \`{a} jour r\'{e}ussissent, soit aucune n'est appliqu\'{e}e.

\begin{figure}[H]
\centering
\begin{tikzpicture}[scale=0.8, transform shape]

% Participants
\node[sdactor] (ui) at (0, 0) {Dashboard};
\node[sdactor] (api) at (4, 0) {NestJS API};
\node[sdactor] (svc) at (8, 0) {StockTransfer};
\node[sdactor] (db) at (12, 0) {PostgreSQL};

% Lifelines
\draw[dashed] (0, -0.35) -- (0, -13);
\draw[dashed] (4, -0.35) -- (4, -13);
\draw[dashed] (8, -0.35) -- (8, -13);
\draw[dashed] (12, -0.35) -- (12, -13);

% Messages
\draw[sdmessage] (0, -1) -- node[above, font=\scriptsize, text width=3cm, align=center] {POST /stock-transfer} (4, -1);
\node[font=\scriptsize\itshape, text=gray] at (2, -1.7) {fromGym, toGym, product, qty};

\draw[sdmessage] (4, -2.5) -- node[above, font=\scriptsize] {createTransfer(dto)} (8, -2.5);

\draw[sdmessage] (8, -3.5) -- node[above, font=\scriptsize, text width=2.5cm, align=center] {SELECT stock source} (12, -3.5);
\draw[sdreturn] (12, -4.5) -- node[above, font=\scriptsize] {qty: 50} (8, -4.5);

\node[draw, dashed, rounded corners, inner sep=3pt, font=\scriptsize] at (10, -5.5) {qty $\geq$ requested?};

\draw[sdmessage] (8, -6.5) -- node[above, font=\scriptsize, text width=2.8cm, align=center] {BEGIN TRANSACTION} (12, -6.5);

\draw[sdmessage] (8, -7.5) -- node[above, font=\scriptsize, text width=2.5cm, align=center] {UPDATE stock source -qty} (12, -7.5);

\draw[sdmessage] (8, -8.5) -- node[above, font=\scriptsize, text width=2.5cm, align=center] {UPDATE stock dest +qty} (12, -8.5);

\draw[sdmessage] (8, -9.5) -- node[above, font=\scriptsize, text width=2.5cm, align=center] {INSERT transfert} (12, -9.5);

\draw[sdmessage] (8, -10.5) -- node[above, font=\scriptsize] {COMMIT} (12, -10.5);

\draw[sdreturn] (8, -11.5) -- node[above, font=\scriptsize] {Transfer cr\'{e}\'{e}} (4, -11.5);
\draw[sdreturn] (4, -12.5) -- node[above, font=\scriptsize] {201 Created} (0, -12.5);

\end{tikzpicture}
\caption{Diagramme de s\'{e}quence -- Transfert de stock inter-salles}
\label{fig:sprint5-sequence}
\end{figure}

\subsection{Interfaces r\'{e}alis\'{e}es (web)}

Pages web : \texttt{articles/} (catalogue), \texttt{entree/} (r\'{e}ception fournisseur), \texttt{sortie/}, \texttt{vente/} (vente membre), \texttt{mouvement-de-stock/} (transferts inter-salles), \texttt{etat-du-stock/} (inventaire filtr\'{e} par cat\'{e}gorie/salle/fournisseur via \texttt{BoutiqueFilter}).

\begin{figure}[H]
  \centering
  \includegraphics[width=\textwidth, height=0.55\textheight, keepaspectratio]{capture/Catalogue des produits.png}
  \caption{Page du catalogue produits avec cards visuelles}
  \label{fig:sprint5-products}
\end{figure}

\begin{figure}[H]
  \centering
  \includegraphics[width=\textwidth, height=0.55\textheight, keepaspectratio]{capture/etat-stock.png}
  \caption{Tableau de bord de l'inventaire par salle de sport}
  \label{fig:sprint5-stock}
\end{figure}

\begin{figure}[H]
  \centering
  \includegraphics[width=\textwidth, height=0.55\textheight, keepaspectratio]{capture/transfert-stock.png}
  \caption{Interface de transfert de stock entre salles}
  \label{fig:sprint5-transfer}
\end{figure}

\begin{figure}[H]
  \centering
  \includegraphics[width=\textwidth, height=0.55\textheight, keepaspectratio]{capture/vente-membres.png}
  \caption{Page de vente de produits aux membres}
  \label{fig:sprint5-sale}
\end{figure}

\begin{figure}[H]
  \centering
  \includegraphics[width=\textwidth, height=0.55\textheight, keepaspectratio]{capture/categories-produits.png}
  \caption{Interface de gestion des cat\'{e}gories de la boutique}
  \label{fig:sprint5-categories}
\end{figure}

\begin{figure}[H]
  \centering
  \includegraphics[width=\textwidth, height=0.55\textheight, keepaspectratio]{capture/Gestion des fournisseurs.png}
  \caption{Page de gestion des fiches fournisseurs}
  \label{fig:sprint5-suppliers}
\end{figure}

\begin{figure}[H]
  \centering
  \includegraphics[width=\textwidth, height=0.55\textheight, keepaspectratio]{capture/reception-stock.png}
  \caption{Interface d'enregistrement des achats et r\'{e}ception de stock}
  \label{fig:sprint5-stock-entry}
\end{figure}

%──────────────────────────────────────────────
\section{Conclusion}
%──────────────────────────────────────────────

De la r\'{e}ception fournisseur \`{a} la vente au membre, la tra\c{c}abilit\'{e} est totale. Le \texttt{StockMovement} horodate chaque mouvement, le transfert transactionnel \'{e}vite les incoh\'{e}rences inter-salles, et le \texttt{ModuleGuard} garantit que seuls les partenaires concern\'{e}s y ont acc\`{e}s. Le Sprint 6 s'attaque \`{a} la couche financière --- facturation, portefeuille et comptabilit\'{e}.

\end{document}
